\section{Conclusion}
\textcolor{black}{In this work, we empirically characterize how agent-system performance varies with coordination structure, model capability, and task properties across 260 controlled configurations spanning three LLM families and six agentic benchmarks. We identify capability saturation as the most robust scaling effect: coordination yields diminishing returns beyond $\sim$45\% single-agent baselines, a finding confirmed under both cluster-robust inference ($p=0.004$) and Holm--Bonferroni multiple-comparison correction ($p_{\text{Holm}}=0.018$). Two additional directional patterns emerge consistently across all six benchmarks: a tool-coordination trade-off where tool-heavy tasks suffer from coordination overhead, and architecture-dependent trace-level error amplification ranging from 4.4$\times$ (centralized) to 17.2$\times$ (independent). Performance gains vary substantially by task structure, from +80.8\% on Finance Agent to $-$70.0\% on PlanCraft, demonstrating that coordination benefits depend on task decomposability rather than team size. We derive a predictive model ($R^2{=}0.373$ across all six benchmarks; $R^2{=}0.413$ with a task-grounded capability metric) that achieves 87\% accuracy in selecting optimal architectures for held-out configurations. On held-out frontier models evaluated on BrowseComp-Plus, the framework remains reasonably calibrated for relative MAS performance prediction (MAS-only MAE=0.061; overall MAE=0.077), providing preliminary evidence of transfer within this benchmark rather than full cross-domain generalization.}